% ===== Introduction. Sources still to add to references.tex / Zotero: =====
% - Afolabi, H., Ram, R. & Rimmel, G. (2022). Harmonization of Sustainability Reporting
%   Regulation: Analysis of a Contested Arena. Sustainability 14(9), 5517. doi:10.3390/su14095517
% - EFRAG (2021). ESRS / double materiality.        - EFRAG (2024). VSME standard.
% - European Commission (2025). VSME recommendation.
% - European Commission (2026). Revised ESRS + value-chain cap (Omnibus). finance.ec.europa.eu
% RESOLVED: The European House - Ambrosetti / TEHA Group (2026) STAGE Index citation
%   (main.tex ~line 260) -- verified real via web search + the link Ruben supplied:
%   https://www.ambrosetti.eu/en/news/sustainable-transition-non-negotiable-priorities-for-the-future-of-businesses/
%   2026 Strategic Report "Non-negotiable priorities for businesses and the future of
%   the sustainable transition" (TEHA Group / The European House-Ambrosetti, with Erion),
%   presented at the Erion Forum, Rome, 17 June 2026. The STAGE Index (Sustainable
%   Transition Alignment between Government & Enterprises) registers ~51% convergence
%   in 2026. In-text now names the index explicitly instead of "one practitioner index".
% RESOLVED: STECF (2025) 3,245-enterprise / two-thirds-micro / large-firm figures --
%   verified via web search against STECF 25-15 "Economic report on the fish processing
%   industry": 3,245 main-activity enterprises in 2023, ~67% microfirms (<10 employees),
%   large enterprises (250+) are ~2% of firms (~65, thesis says 66 -- a one-firm rounding
%   difference from the 2%-of-3,245 back-calculation; worth a 5-minute check against the
%   report's own table if you want the exact reported count rather than the recomputed one).
% RESOLVED: the "alphabet soup / harmonisation" cite at main.tex:414 now points to
%   Afolabi, Ram & Rimmel (2022), whose title is literally about this contested
%   harmonisation arena -- a tighter match than Adams & Abhayawansa. Verified real
%   via search: Sustainability 14(9), 5517, doi:10.3390/su14095517.
% RESOLVED: the "buyers drop suppliers over ESG data" line in 1.3 now cites
%   Tanase (2025), Veridion -- a real, dated industry source confirming ESG
%   screening/scoring gates suppliers out at exactly the point described.

\section{1. Introduction}\label{introduction}

\subsection{1.1 A crowded and fragmented field of
demands}\label{a-rising-tide-of-sustainability-demands}

Companies today face a large and steadily growing volume of requests for
sustainability information, and from a widening set of actors. Regulators
legislate mandatory disclosures. Retail and business customers send
sustainability questionnaires and make certification a condition of supply.
Banks, investors and ratings agencies ask for climate and ESG data before
extending capital or assigning a score. Industry coalitions and NGOs promote
voluntary standards and commitments. Each actor asks through its own channel, on
its own timetable, and in its own vocabulary.

These demands are less independent than they first appear, because many of them
share a common origin. The European Union has made corporate sustainability
reporting a central lever of its Green Deal: the Corporate Sustainability
Reporting Directive (CSRD), implemented through the European Sustainability
Reporting Standards (ESRS), defines mandatory disclosures across environmental,
social and governance topics on a double-materiality basis (EFRAG, 2021). These
obligations bind only large undertakings directly, but they do not stop there. A
company in scope must gather data from its suppliers, and passes the requirement
down its value chain as contractual questionnaires. Banks and investors,
themselves regulated under the SFDR, ask the same of the companies they finance.
A single regulatory requirement therefore resurfaces, in different forms,
through several actors at once. What a small supplier experiences as 'a
customer demand' is frequently a regulatory demand that has cascaded one or two
steps down the chain (a pattern this thesis establishes directly in Chapter~4's
source-coding test, Section~4.1).
This is a dynamic the EU has itself acknowledged by
legislating, in its 2026 simplification package, a 'value-chain cap' to limit
how far such requests may be pushed onto small suppliers (European Commission,
2026).

Seen one request at a time, the combined effect is a proliferating,
uncoordinated wall of separate demands: the 'alphabet soup' of standards and
frameworks that has prompted sustained calls for harmonisation (Afolabi, Ram \&
Rimmel, 2022).
Goals are not the sticking point so much as alignment between the actors pursuing
them: the STAGE Index, which measures how far government and enterprise
sustainability agendas coincide, registers only about 51 per cent convergence in
2026 (The European House~--~Ambrosetti, 2026). For the firm on the receiving end,
the practical experience is simply that the volume and variety of what is asked
keeps growing, and that keeping up takes ever more effort.

\subsection{1.2 Different questions, shared
answers}\label{different-questions-shared-answers}

The observation motivating this thesis is that, beneath their different wording,
these demands draw on a much smaller body of shared data. A customer's
carbon-footprint question, an investor's climate metric and a regulatory
emissions disclosure all rest on the same greenhouse-gas accounting. A
retailer's traceability requirement, a certification's chain-of-custody audit, a
due-diligence law and a value-chain disclosure standard all rest on the same
lot-level tracking data. Behind a large number of superficially different and
fragmented questions sits a much smaller number of shared datapoints, and those
datapoints cluster into a handful of recurring themes. Establishing that this
overlap is real, and measuring how large it is, is precisely the work of the
analysis that follows.

That overlap is not evenly spread. Some datapoints are demanded by nearly every
actor, others by only one. A few central themes are where the most actors
converge, and these underpin a firm's licence to operate and to sell; a long tail
of remaining topics is asked less often and by fewer actors.
The practical implication is direct. A firm that identifies the
key themes and learns to produce their datapoints \emph{once} - as a
structured, reusable process rather than a one-off answer - can satisfy most of
what it is asked from a single effort, instead of rebuilding an answer for every
request.

\subsection{1.3 Capacity, not intent: why small firms stay
reactive}\label{why-resource-constrained-firms-stay-reactive}

The barrier to capturing this prize - answering many demands from one
well-built set of data -- is not a matter of intent but of capacity. Small and
medium-sized enterprises (SMEs) operate with limited financial and managerial
resources, little dedicated reporting expertise and weak data infrastructure.
SMEs do not operate simply as miniature large firms, and their performance
should be assessed against benchmarks that fit the circumstances of small
firms rather than standards built for large-firm capacity (Battisti \&
Perry, 2011).
Under those constraints the up-front investment in
structured data processes, the very investment that would let a firm answer
many frameworks from one dataset, is rarely made. Instead the firm responds
\emph{reactively}: each new questionnaire, audit or regulation is handled as a
separate, one-off exercise, its data assembled by hand and then set aside
(the pattern the scoping conversations described in Section~1.4 pointed to).
Fragmented governance thus pushes resource-constrained firms toward
minimum-viable compliance rather than strategic integration (Chan, Lai \& Kim,
2022).

The consequence follows directly: because little or no reusable capability was
built, each new regulation or customer request has to be met from a standing
start. Data is assembled again for each request, and the firm stays a step
behind - mostly catching up to the last demand, rarely prepared for the next - the reactive posture developed fully in
Chapter~2 (Section~2.2). Two licences are at stake at once, the
\emph{legal} licence to operate, which direct regulation makes non-negotiable,
and the \emph{commercial} licence to sell, which large customers and their
financiers increasingly attach to the same sustainability data.
The stakes are concrete. Buyers increasingly require suppliers to complete ESG
assessments before contracting. A supplier that cannot answer risks being scored
down, dropped from the supply chain, or screened out at the rating stage before a
relationship starts. Both are documented dynamics of buyer-side screening
platforms such as EcoVadis (Tanase, 2025), whose reach across supplier networks
makes the assessment a de facto condition of trade rather than an optional
exercise.
A firm that builds the underlying data once avoids repeating that risk with
every buyer; one that answers each demand in isolation carries it every time.

\subsection{1.4 Seafood processing in Europe: a revealing
case}\label{seafood-processing-a-revealing-case}

To make the problem tangible and empirically grounded, this thesis studies one
sector in depth: European seafood processing. Three features make it a revealing
case.

First, it is overwhelmingly a sector of small firms. Around two-thirds of the
roughly 3{,}245 EU enterprises that process fish as their main activity are
micro-enterprises of ten employees or fewer, and only about two per cent employ
250 people or more (STECF, 2025). The great majority of the sector therefore sits
in the resource-constrained position described above.

Second, it carries a dense and fragmented governance load. Seafood processors
sit beneath the full stack of horizontal EU sustainability law (the CSRD, CSDDD,
the Packaging and Packaging Waste Regulation, the Empowering Consumers Directive
and more), a layer of sector-specific regulation (the Common Fisheries Policy,
the IUU and Control Regulations, the EU Deforestation Regulation), \emph{and} a
voluntary layer that stacks certification (MSC and RFM in wild capture; ASC,
GLOBALG.A.P.\ and BAP in aquaculture), buyer rating platforms (EcoVadis, Sedex)
and traceability standards (GDST) on top of one another. A fourth strand runs
across all three: ethical and human-rights demand, reaching the firm through the
CSDDD's due-diligence duty, the Forced Labour Regulation and buyer codes of
conduct. Few of these are
individually onerous - and several are interoperable, with MSC and ASC both
working with GDST. The load comes from their number as a \emph{stack}: each is
answered separately, on its own timetable and in its own format.

Third, this compounding pressure is not merely inferred from the regulatory map.
Informal scoping conversations with small seafood processors in Belgium and
Germany, held at the outset of this research, pointed the same way. Those firms
described simultaneous, overlapping requests from regulators, retail customers
and finance providers, which they answered one at a time because they had no
shared process for answering them together.\footnote{These conversations were
exploratory and unstructured, conducted to orient the research rather than as a
data-collection method. They are not coded, quoted or treated as evidence
anywhere in this thesis; the empirical claims rest entirely on the document
analysis of Chapters~4 to~7.}

Seafood is a case, however, not the point. The same structure, many actors,
overlapping underlying data, resource-constrained firms, recurs across
manufacturing and agri-food SME sectors facing the same widening set of
demands. Seafood processing is a sharp instance that makes the general
problem visible and testable; the diagnostic approach developed here is
intended to be transferable.

\subsection{1.5 Aim, approach and
contribution}\label{aim-approach-and-contribution}

The EU has not been blind to this burden. Its main simplification response -
the Voluntary Sustainability Reporting Standard for SMEs (VSME), which EFRAG
developed and the Commission endorsed by recommendation in 2025, alongside the
wider Omnibus revision now cutting ESRS datapoints by more than seventy per cent
(European Commission, 2026) - lightens the \emph{volume} of what a small firm
must report. But it standardises \emph{what} a proportionate disclosure
contains, not \emph{which} themes an SME should build first amid its own
particular mix of regulators, customers and financiers. It answers what a
lighter report looks like; it does not answer where, among those demands, a
firm should begin.

This thesis addresses that second question: where, amid the overlapping
demands, a resource-constrained European seafood-processing SME should begin.
Its aim is to identify the key disclosure themes - and the underlying
datapoints - that, if built first, satisfy the largest share of the demands
the firm faces across regulators, customers, investors and industry initiatives,
and thereby to support a shift from reactive compliance to proactive strategy
while helping secure both the licence to operate and the licence to sell. Neither
licence is new: buyers have long imposed financial, quality and food-safety
conditions beyond what the law requires, and in some areas - microbiological
limits, or human-rights due diligence - the legal floor is the stricter of the
two. What is new is the volume and formalisation of the sustainability-specific
demand now attached to both.

It approaches the problem in four steps, each forming one empirical chapter.
Step~1 (Chapter~4) scans the \emph{ecosystem}: the full landscape of
instruments acting on the sector, who demands what, organised by how
obligatory each demand is and which stakeholder drives it, using the
mandatory ESRS as a common reference frame. Step~2 (Chapter~5) scans the
\emph{peers}: benchmarked against that ESRS structure, what some of the
sector's most advanced reporters actually disclose, establishing the
realistic frontier of seafood disclosure and a universal core of topics
every mature firm reports. Step~3 (Chapter~6) \emph{crosses the two}: mapping
which instruments demand which topics reveals where demand converges, the
hotspots at which a single theme, once its datapoints are in place, answers
many actors at once. Step~4, carried out in the same chapter because staging
and tool-building draw on the same underlying matrix, \emph{builds the
tool}: it stages the surviving topics by convergence and stakeholder breadth
and operationalises the result as a diagnostic with which an individual SME
can locate its own position and identify what to build first.
Figure~\ref{fig:methodology-overview} summarises the four steps. Chapter~3
states the research questions and intended contributions; Chapter~2 reviews
the literature and develops the conceptual framework (the reactive-proactive
distinction, stakeholder pressure and legitimacy); Chapter~3 sets out the
shared research design. Chapter~7 draws out the implications for seafood
SMEs and for EU disclosure policy, and concludes.

The contribution is threefold. For seafood SMEs, the thesis offers a diagnostic
framework and a prioritised, sequenced set of topics tailored to a firm's
situation, intended as a basis for strategic discussion within the firm rather
than a mere compliance checklist. For the academic literature, it provides a
systematic mapping of cross-stakeholder disclosure demand at sector level and an
empirically grounded application of the reactive/proactive distinction. For EU
policy, it supplies evidence to the ongoing simplification debate, the Omnibus
package, the VSME, the postponement of sector standards, by identifying where
the disclosure ecosystem converges, and therefore where coordinating or capping
demand would deliver the largest reductions in SME burden. The analysis that
follows shows precisely which themes those are and which datapoints they
require.
